% =====================================================================
% PLANTILLA DE ENTREGA EN PDF (estilo aprobado: simulacros/simulacro-01)
% Compilar: pdflatex plantilla-pdf.tex   (dos veces)
% =====================================================================
\documentclass[11pt]{article}
\usepackage[spanish,es-nodecimaldot]{babel}   % punto decimal, como el profesor
\usepackage[utf8]{inputenc}
\usepackage{amsmath,amssymb,geometry,array,tikz}
\usetikzlibrary{babel}   % evita conflicto de babel-spanish con flechas de TikZ
\geometry{letterpaper,margin=1.5cm}
\AtBeginDocument{\setlength{\abovedisplayskip}{3pt}\setlength{\belowdisplayskip}{3pt}\setlength{\abovedisplayshortskip}{2pt}\setlength{\belowdisplayshortskip}{2pt}}
\setlength{\parindent}{0pt}
\renewcommand{\arraystretch}{1.3}
\DeclareMathOperator{\sen}{sen}
\begin{document}
\begin{center}{\Large\textbf{TÍTULO --- Matemáticas Aplicadas}}\end{center}
\hrule\vspace{8pt}

% ---------- Bloque de ejercicio (uno por página: \newpage entre ejercicios) ----------
{\large\textbf{1) Título del tema como lo escribe el profesor}}\\[6pt]
\textbf{Resolver} \quad $\text{ecuación} \qquad \text{condiciones} \qquad H=0.1$\\[4pt]

% Sistema f/g (R-K):
\[
\begin{array}{r} f:\\[2pt] g:\end{array}\left\{\begin{array}{l} y'=z\\[2pt] z'=\dots\end{array}\right.
\qquad\qquad y(t_0)=\dots \qquad z(t_0)=\dots
\]

% Pasos: align* con un paso por línea, sustitución completa -> resultado (4 decimales truncados)
\begin{align*}
K_1&=0.1f(\dots)=0.1[\dots]=\dots\\
C_1&=0.1g(\,''\,)=0.1[\dots]=\dots
\end{align*}

% Resultado final recuadrado (tabla o valor)
\[\boxed{\begin{array}{c|ccc} t & \dots \\\hline y & \dots \end{array}}\]

% Nota mínima solo si hace falta (supuestos o errores del enunciado)
{\footnotesize $^{*}$ Nota breve.}
\end{document}
